\documentclass[12pt]{article}
\usepackage[utf8]{inputenc}
\usepackage[T1]{fontenc}
\usepackage{hyperref}
\title{Warrwa Language Resources for K-12 Curriculum Development}
\author{Compiled by Claude Sonnet 5.5 (claude-sonnet-5-5) for Andr\'as Kornai}
\date{}
\begin{document}
\maketitle

\section{Introduction}
Warrwa (ISO 639-3 wwr; Glottocode warr1258) is an extinct Nyulnyulan language of the Kimberley region of Western Australia. It belongs to the Eastern Nyulnyulan group, together with Nyikina, Yawuru and Jukun. Search results report that the last known speaker, Maudie Lennard, died in 2016. Some sources treat Warrwa as a dialect of Nyigina, and others treat it as a separate language. Very little teaching material appears in the search results, so any curriculum would be a revitalization effort built on archival and descriptive sources.

\section{Primary Educational Resources}
No primers, textbooks or school programmes for Warrwa appeared in the searches. None are claimed here.

\section{Linguistic and Reference Materials}
\begin{itemize}
\item William McGregor, \emph{A Grammar of Warrwa, Kimberley, Western Australia} (Pacific Linguistics, De Gruyter). It is based on fieldwork with the last speakers, Maudie Lennard and Freddy Marker. It includes a wordlist and texts. \url{https://www.degruyterbrill.com/document/isbn/9783111401959/html}
\item AIATSIS AUSTLANG entry K10 for Warrwa, with language and territory information. \url{https://collection.aiatsis.gov.au/austlang/language/K10}
\item Grambank entry for Warrwa (warr1258). \url{https://grambank.clld.org/languages/warr1258}
\item Wikipedia articles on the Warrwa language and on the Nyulnyulan family, which includes Capell's 1940 comparative vocabulary. \url{https://en.wikipedia.org/wiki/Warrwa_language} and \url{https://en.wikipedia.org/wiki/Nyulnyulan_languages}
\end{itemize}

\section{Religious and Community Resources}
No Warrwa scripture translations or community organisations were found in the searches.

\section{Digital and Audio Resources}
No Warrwa audio, video or software was found in the searches. McGregor's fieldwork materials may exist in archives, but I saw no URL for them.

\section{Curriculum Development Context}
The language is extinct, so the realistic aim is a heritage and revitalization curriculum. It could draw on McGregor's grammar, wordlist and texts, and on the AIATSIS records. Any use should involve the Warrwa and wider Kimberley communities and their descendants. The first step is to contact the relevant Aboriginal language organisations. None were identified in the searches.
\end{document}